\chapter{Drafts}
\section{Kotlin Multiplatform}
\label{sec:kotlin-multiplatform}
One of the main advantages of Kotlin is its support for multiplatform programming.
It helps keeping the flexibility and advantages of native programming while cutting down on
the effort spent writing and maintaining the same code for several platforms.

\centerImage{
	figures/kotlin-multiplatform.png
}{
	Kotlin Multiplatform - Supported Platforms
}{
	kotlin-mp
}{
	0.5
}

\begin{itemize}
	\item The language, basic tools, and core libraries are all part of \emph{Common Kotlin}. Common Kotlin code runs on all systems and is portable.
	\item It is possible to use Kotlin platform-specific versions for platform interoperability. Extensions to the Kotlin language as well as platform-specific libraries and tools are included in the \emph{Kotlin/JVM}, \emph{Kotlin/JS}, and \emph{Kotlin/Native} platform-specific versions of the language.
	Through these platforms it is possible to access the platform native code (JVM, JS, and Native) and utilizing all the native features.
\end{itemize}

Kotlin Multiplatform may have many applications, but the following are some of the most prevalent ones:

\begin{itemize}
	\item Create Android, iOS, and web apps. By using shared code, the amount of business logic that frontend developers have to create is decreased.
	This improves the efficiency of product implementation while requiring less coding and testing for the specific platform.
	\item Create of a multiplatform library.
	Using shared code and platform-specific components if needed, a multiplatform library that has been distributed may become a classic dependency.
	\item Building a full-stack web app. The two main targets are Kotlin/JVM (for the server) and Kotlin/JS (for the client). The shared code, mainly constints in the model of the application, but can also contains common logic that may run on both the server and the client.
\end{itemize}
